\documentclass{ximera}
\author{Jeffrey Kuan}
\title{Math 1050 HW 5.1 --- The Rectangular Coordinate System}
\license{CC: 0}

\begin{document}

\begin{abstract}
    Section 5.1 --- The Rectangular Coordinate System.
    Assigned exercises: 1, 3, 5, 6, 11–19 odd, 25, 33, 34, 37, 39, 40.
\end{abstract}

\maketitle

SECTION 5.1 1, 3, 5, 6, 11 –- 19 odd, 25, 33, 34, 37, 39, 40.

\begin{problem}
\textbf{1.} In which quadrant is the graph of $(-3, 4)$?

\begin{multipleChoice}
\choice{First quadrant}
\choice[correct]{Second quadrant}
\choice{Third quadrant}
\choice{Fourth quadrant}
\end{multipleChoice}
\end{problem}

\begin{problem}
\textbf{3.} On which axis is the graph of $(0, -4)$?

\begin{multipleChoice}
\choice{x-axis}
\choice[correct]{y-axis}
\end{multipleChoice}
\end{problem}

\begin{problem}
\textbf{5.} What is the value of the $y$-coordinate of any point on the $x$-axis?

$\answer{0}$
\end{problem}

\begin{problem}
\textbf{6.} What is the value of the $x$-coordinate of any point on the $y$-axis?

$\answer{0}$
\end{problem}

\begin{problem}
\textbf{11.} Graph the ordered pairs $(-2, 1)$, $(3, -5)$, $(-2, 4)$, and $(0, 3)$.
  \figbox{\gridsm}
\end{problem}

\begin{problem}
\textbf{13.} Graph the ordered pairs $(0, 0)$, $(0, -5)$, $(-3, 0)$, and $(0, 2)$.
  \figbox{\gridsm}
\end{problem}

\begin{problem}
\textbf{15.} Graph the ordered pairs $(-1, 4)$, $(-2, -3)$, $(0, 2)$, and $(4, 0)$.
  \figbox{\gridsm}
\end{problem}

\begin{problem}
\textbf{17.} Find the coordinates of each of the points.
  \figbox{\gridpts{%
    \pt{2}{3}{A}{above left}
    \pt{4}{0}{B}{above left}
    \pt{-4}{1}{C}{above right}
    \pt{-2}{-2}{D}{below left}}}

Enter your answer as an ordered pair with no spaces.

 Point A has coordinates $ \answer{(2,3)} $ 

 Point B has coordinates $ \answer{(4,0)} $ 

 Point C has coordinates $ \answer{(-4,1)} $ 

 Point D has coordinates $ \answer{(-2,-2)} $
\end{problem}

\begin{problem}
\textbf{19.} Find the coordinates of each of the points.
  \figbox{\gridpts{%
    \pt{-2}{5}{A}{below right}
    \pt{3}{4}{B}{below right}
    \pt{0}{0}{C}{above right}
    \pt{-3}{-2}{D}{below left}}}

Enter your answer as an ordered pair with no spaces.

Point A has coordinates $\answer{(-2,5)}$.

Point B has coordinates $\answer{(3,4)}$.

Point C has coordinates $\answer{(0,0)}$.

Point D has coordinates $\answer{(-3,-2)}$
\end{problem}

\begin{problem}
\textbf{25.} \textbf{Fuel Efficiency} The American Council for an Energy-Efficient Economy
  releases rankings of environmentally friendly and unfriendly cars and trucks sold in the
  United States. The following table shows the fuel usage, in miles per gallon, in the city
  and on the highway, for the twelve 2010-model vehicles ranked best for the environment.
  Graph the scatter diagram for these data.
  \figbox{%
  \begin{tabular}{@{}l|cccccccccccc@{}}
  \toprule
  Fuel usage in mpg, city, $x$    & 51 & 40 & 29 & 28 & 34 & 22 & 22 & 22 & 21 & 18 & 21 & 21 \\
  Fuel usage in mpg, highway, $y$ & 48 & 45 & 35 & 35 & 31 & 32 & 31 & 32 & 27 & 25 & 22 & 22 \\
  \bottomrule
  \end{tabular}}
  \figbox{%
  \begin{tikzpicture}[x=0.10cm,y=0.10cm]
    \draw[gridln] (18,18) grid[step=2] (54,52);
    \draw[axisln,-{Stealth[length=4pt]}] (18,18) -- (56,18) node[gl,right] {$x$};
    \draw[axisln,-{Stealth[length=4pt]}] (18,18) -- (18,54) node[gl,above] {$y$};
    \foreach \v in {20,30,40,50} {\node[gl,below] at (\v,17) {\v};}
    \foreach \v in {20,30,40} {\node[gl,left] at (17,\v) {\v};}
    \node[gl,below=8pt,align=center] at (37,17)
      {Fuel usage, city\\(in miles per gallon)};
    \node[gl,rotate=90,align=center] at (8,36)
      {Fuel usage, highway\\(in miles per gallon)};
  \end{tikzpicture}}
\end{problem}

\begin{problem}
\textbf{33.} \textbf{Temperature} On September 10 in midstate New Hampshire, the temperature at
  6 \textsc{a.m.} was $45^\circ$F. At 2 \textsc{p.m.} on the same day, the temperature was
  $77^\circ$F. Find the average rate of change in temperature per hour.

The average rate of change is $\answer{4}$ degrees Fahrenheit per hour.

34(a) $\answer{11204}$ people per year

(b) \begin{multipleChoice}
\choice[correct]{greater than}
\choice{less than}
\end{multipleChoice}

(c) During the decade starting in the year $\answer{1610}$, the average rate of change was $\answer{195}$ people per year. 

37(a) $\answer{26}$ students per year.

(b)$\answer{566}$ students per year.

(c) \begin{multipleChoice}
\choice{between 2000 and 2005}
\choice[correct]{between 2005 and 2009}
\end{multipleChoice}
\end{problem}

\begin{problem}
\textbf{34.} \textbf{Demography} The table below shows the population of Colonial America for
  each decade from 1610 to 1780. \emph{(Source: infoplease.com)} These figures are
  estimates, because this period was prior to the establishment of the U.S. Census in 1790.
  \begin{enumerate}[label=\textbf{\alph*.},leftmargin=2em,topsep=2pt,itemsep=1pt]
    \item Find the average annual rate of change in the population of the Colonies from
          1650 to 1750.
    \item Was the average annual rate of change in the population from 1650 to 1750
          greater than or less than the average annual rate of change in the population
          from 1700 to 1750?
    \item During which decade was the average annual rate of change the least? What was
          the average annual rate of change during that decade?
  \end{enumerate}
  \figbox{%
  \footnotesize
  \setlength{\tabcolsep}{4pt}
  \begin{tabular}{@{}l|ccccccccc@{}}
  \toprule
  Year       & 1610 & 1620 & 1630 & 1640 & 1650 & 1660 & 1670 & 1680 & 1690 \\
  Population & 350 & 2,300 & 4,600 & 26,600 & 50,400 & 75,100 & 111,900 & 151,500 & 210,400 \\
  \midrule
  Year       & 1700 & 1710 & 1720 & 1730 & 1740 & 1750 & 1760 & 1770 & 1780 \\
  Population & 250,900 & 331,700 & 466,200 & 629,400 & 905,600 & 1,170,800 & 1,593,600 & 2,148,100 & 2,780,400 \\
  \bottomrule
  \end{tabular}}
\end{problem}

\begin{problem}
\textbf{37.} \textbf{Education} Read the news clipping below.
  \newsbox{More Students Opt to Be Cool with Summer School}{%
    The advantage of lower costs and the chances to get ahead in their coursework lures an
    increasing number of students at colleges and universities to choose summer sessions.
    Here's a look at the summer enrollment numbers for two institutions.
    \emph{Source: The Boston Globe, June 22, 2010}}
  \figbox{%
  \small
  \begin{tabular}{@{}l|ccc@{}}
  \multicolumn{4}{@{}l}{\textbf{Boston College Summer Enrollment}}\\
  \toprule
  Year & 2000 & 2005 & 2009 \\
  Number of students enrolled & 3742 & 3873 & 6137 \\
  \bottomrule
  \end{tabular}}
  \figbox{%
  \small
  \begin{tabular}{@{}l|ccc@{}}
  \multicolumn{4}{@{}l}{\textbf{University of Massachusetts--Amherst Summer Enrollment}}\\
  \toprule
  Year & 2000 & 2005 & 2009 \\
  Number of students enrolled & 6254 & 7980 & 9576 \\
  \bottomrule
  \end{tabular}}
  \begin{enumerate}[label=\textbf{\alph*.},leftmargin=2em,topsep=2pt,itemsep=1pt]
    \item What was the average annual rate of change in summer enrollment at Boston
          College between 2000 and 2005? Round to the nearest whole number.
    \item What was the average annual rate of change in summer enrollment at Boston
          College between 2005 and 2009? Round to the nearest whole number.
    \item Which was greater at the University of Massachusetts--Amherst: the average
          annual rate of change in summer enrollment between 2000 and 2005, or the average
          annual rate of change in summer enrollment between 2005 and 2009?
  \end{enumerate}
\end{problem}

\begin{problem}
\textbf{39.} \textbf{Minimum Wage} The following table shows buying power of a worker earning
  the federal minimum wage for various years. This is not the \emph{actual} minimum wage
  but the minimum wage adjusted for inflation. For instance, the actual minimum wage in
  1970 was \$1.45. Spending \$1.45 in 1970 was the same as spending \$8.12 in 2010.
  \emph{(Source: en.wikipedia.org)}
  \figbox{%
  \small
  \begin{tabular}{@{}l|cccccccc@{}}
  \toprule
  Year & 1970 & 1975 & 1980 & 1981 & 1990 & 1996 & 2007 & 2009 \\
  Minimum wage & 8.12 & 8.08 & 8.18 & 8.01 & 6.33 & 6.59 & 6.14 & 7.36 \\
  \bottomrule
  \end{tabular}}
  \begin{enumerate}[label=\textbf{\alph*.},leftmargin=2em,topsep=2pt,itemsep=1pt]
    \item What was the average annual rate of change in the inflation-adjusted minimum
          wage from 1970 to 1990? \emph{(Note: A negative average rate of change denotes a
          decrease.)}
    \item Prior to 1981, in which five-year period did the average annual rate of change
          in the inflation-adjusted minimum wage increase? What was the average annual rate
          of change during that period?
  \end{enumerate}

(a) Rounded to the nearest cent, the average rate of change was $\answer{-9}$ cents per year.

(b) From the year $\answer{1975}$ to $\answer{1980}$ the rate of change was $\answer{2}$ cents per year.

40(a) Rounded to the nearest million, the average rate of change was $\answer{-2}$ million per year.

(b) The average rate of change was (enter your answer without commas) $\answer{-1300000}$ people per year.
\end{problem}

\begin{problem}
\textbf{40.} \textbf{Demography} The following table shows the past and projected populations of
  baby boomers for the years 2000, 2031, and 2046. \emph{(Source: U.S. Census Bureau)}
  \figbox{%
  \small
  \begin{tabular}{@{}l|ccc@{}}
  \toprule
  Year & 2000 & 2031 & 2046 \\
  Population, in millions & 79 & 51 & 19 \\
  \bottomrule
  \end{tabular}}
  \begin{enumerate}[label=\textbf{\alph*.},leftmargin=2em,topsep=2pt,itemsep=1pt]
    \item Find the average rate of change per year in the population of the baby boomer
          generation from 2031 to 2046. \emph{(Note: A negative average rate of change
          denotes a decrease in population.)}
    \item Find the average rate of change per year in the population of baby boomers from
          2000 to 2046. Round to the nearest hundred thousand.
  \end{enumerate}
\end{problem}

\end{document}
